% Section 7 — Case studies (NEW for GMD). Each case study uses
% sbatch_case_study.sh + fig_case_study.py. Numbers TBD until SLURM
% jobs run, but the narrative structure is in place.

We document three case studies of operational interest from 2022.
Each case is reproducible end-to-end through a single SLURM
submission (\texttt{paper/scripts/sbatch\_case\_study.sh}) followed
by the figure script (\texttt{paper/scripts/fig\_case\_study.py}).
For each case we show four-panel maps (observed, CRAN-PM, CAMS
forecast, persistence) on the day of peak intensity, and station
time series at six representative EEA monitors.

\subsection{Saharan dust intrusion (14--19 March 2022)}
\label{sec:case:saharan}

A massive Saharan dust outbreak ("Calima") affected southern and
western Europe in mid-March 2022, with PM\textsubscript{10} and
PM\textsubscript{2.5} concentrations exceeding the EU 24-hour
limit value at hundreds of stations from Spain to the United
Kingdom. The episode is well documented in remote-sensing imagery
(MODIS aerosol optical depth) and required interventions from
several national air-quality agencies.

This case stresses the model's ability to reproduce
\emph{long-range} aerosol transport, which is mediated by
synoptic-scale wind patterns. The CAMS regional forecast captures
the broad outlines of the episode but underestimates the peak
near-surface concentrations because of its coarser resolution and
its limited representation of the local boundary-layer dynamics.

% \begin{figure}[!ht]
%   \centering
%   \includegraphics[width=\textwidth]{figures_cases/saharan-dust/fig_case_saharan_dust_maps.pdf}
%   \caption{\textbf{Saharan dust intrusion, March 2022 case study.}
%   Four-panel comparison on the peak day. Top row: observed (GHAP),
%   CRAN-PM forecast. Bottom row: operational CAMS forecast, persistence.}
%   \label{fig:case_saharan}
% \end{figure}

% TODO: insert maps and station time series after running
% `sbatch sbatch_case_study.sh saharan-dust 2022-03-14 2022-03-19`.

\subsection{Iberian Peninsula wildfires (14--22 July 2022)}
\label{sec:case:iberian}

Summer 2022 was the worst wildfire season on record in Spain and
Portugal. The Sierra de la Culebra fire (Spain) produced large
PM\textsubscript{2.5} plumes that travelled hundreds of kilometres.

This case stresses the model's ability to reproduce
\emph{transient, locally concentrated} emissions of biomass-burning
aerosol, which are challenging for chemistry-transport models that
rely on satellite-derived emission inventories with several days of
latency. Because CRAN-PM ingests the GHAP PM\textsubscript{2.5}
analysis at $t-1$, it inherits the satellite plume signature and
propagates it forward.

% TODO: insert maps and station time series after running
% `sbatch sbatch_case_study.sh iberian-fires 2022-07-14 2022-07-22`.

\subsection{Polish/Slovak winter heating episode (8--18 January 2022)}
\label{sec:case:polish}

A cold snap in early January 2022 drove residential coal and biomass
combustion to peak levels in Central Europe. Krakow, Katowice
(Poland), Ostrava (Czech Republic) and Bratislava (Slovakia) all
recorded multi-day PM\textsubscript{2.5} concentrations above
75\,$\mu$g\,m$^{-3}$, four times the EU annual limit.

This case stresses the model's ability to resolve
\emph{urban-scale gradients} (within a $\sim 50$\,km radius) that
the operational CAMS forecast smears out. The case is also
politically important: persistent winter exceedances in Central
Europe are central to the EU Air Quality Directive infringement
proceedings and to the funding of clean-heating subsidy programmes.

% TODO: insert maps and station time series after running
% `sbatch sbatch_case_study.sh polish-winter 2022-01-08 2022-01-18`.
